\section{Introduction}\label{sec:intro}

Chronic kidney disease (CKD) affects over 1 in 7 American adults and can progress to end-stage renal disease (ESRD) \citep{CDC2023CKD}. Early identification of high-risk patients is critical for nephrology referral and dialysis-access planning \citep{levin2011early_detection}. The clinical standard for this task is the Kidney Failure Risk Equation (KFRE) \citep{tangri2011predicting}, a closed-form Cox model built from as few as four routinely collected variables (age, sex, eGFR, urine albumin-to-creatinine ratio (uACR)), or eight in its extended form (adding calcium, phosphate, bicarbonate, albumin). KFRE has been externally validated across 31 cohorts spanning over 700{,}000 patients \citep{tangri2016multinational} and is embedded in clinical guidelines \citep{major2019kfre_external_validation}.

Modern EHR resources and flexible survival models motivate comparison with KFRE across different input representations. Adding laboratory variables can also change cohort eligibility and the available history, complicating interpretation of feature-count comparisons.

We benchmark ten learned survival models in three MIMIC-IV-derived scenarios: \textbf{four\_features} and \textbf{eight\_features} match KFRE's four- and eight-variable inputs; \textbf{twenty\_features} uses 20 laboratory variables with missingness flags and has no published KFRE analogue. We add the closed-form KFRE equation to the first two scenarios (11 models there). The extracted cohorts contain 2{,}809, 1{,}213, and 32{,}601 patients, respectively. They share a source population but do not form a controlled, nested feature-count experiment.

Building complete KFRE-input rows requires merging asynchronous measurements around creatinine anchors while retaining repeated observations (Section~\ref{sec:cohort}). Nearest-value matching can include measurements after an anchor, so this design remains retrospective.

\textbf{Contributions.} We provide a documented cohort-construction procedure and a completed comparison across all three scenarios, reporting mean and sample SD over five runs on the same patient partition. We use one prediction per patient while distinguishing final-row and full-history inputs. We also document the limits of the completed analysis: scenario-specific cohorts, test-informed LogisticHazard selection, raw-row censoring references, and unmatched decision-curve comparators. The results describe these implementations and do not establish the isolated effect of feature count or clinical utility.
